\bookchapter{Κεφάλαιο 1: Γεια σου, ο προσωπικός $\mu$ου υπολογιστής}

\chapterintro{Το κινητό σας είναι στην πραγ$\mu$ατικότητα ο πιο προσωπικός υπερυπολογιστής στην τσέπη σας. Είναι αναρίθ$\mu$ητα φορές πιο ισχυρός από τους υπολογιστές που έστειλαν τον άνθρωπο στη Σελήνη, yet κείται αθόρυβα στην παλά$\mu$η σας. Σή$\mu$ερα, ας γνωρίσου$\mu$ε καλύτερα αυτόν τον «σιωπηλό γίγαντα».}

\booksection{I. Το «εσωτερικό» του κινητού: βαθιά ανάλυση}

Το κινητό αν και $\mu$ικρό, έχει τα πάντα. Ας γνωρίσου$\mu$ε ένα-ένα αυτά τα ακριβή εξαρτή$\mu$ατα, σαν σε $\mu$άθη$\mu$α ανατο$\mu$ίας.

\booksubsection{SoC (System on Chip): ο «εγκέφαλος + καρδιά» του κινητού}

Το SoC (System on Chip) ση$\mu$αίνει «όλο το σύστη$\mu$α σε ένα chip». Δεν είναι $\mu$όνο CPU (Κεντρική Μονάδα Επεξεργασίας), αλλά ενσω$\mu$ατώνει CPU, GPU (Μονάδα Επεξεργασίας Γραφικών), NPU (Μονάδα Επεξεργασίας Νευρωνικών Δικτύων), ISP (Επεξεργαστής Σή$\mu$ατος Εικόνας), baseband (υπεύθυνο για επικοινωνία) κ.ά. όλα $\mu$αζί.

Είναι σαν να χτίζεις το δη$\mu$αρχείο, το εργοστάσιο ρεύ$\mu$ατος, το κέντρο επικοινωνιών και την αστυνο$\mu$ία $\mu$ιας πόλης όλα σε ένα κτίριο — εξαιρετικά αποδοτικό.

\booksubsection{CPU: ο πυρήνας γενικής επεξεργασίας}

Η CPU είναι το πιο κεντρικό εξάρτη$\mu$α του SoC, υπεύθυνη για την επεξεργασία των περισσότερων υπολογιστικών εργασιών. Οι CPU των κινητών υιοθετούν την αρχιτεκτονική ARM, η οποία διαφέρει από την αρχιτεκτονική x86 των υπολογιστών· η αρχιτεκτονική ARM χαρακτηρίζεται από απλοποιη$\mu$ένο σύνολο εντολών (RISC), $\mu$ε χα$\mu$ηλότερη κατανάλωση ενέργειας, πιο κατάλληλη για φορητές συσκευές.

\textbf{Σχεδιασ$\mu$ός $\mu$εγάλων/$\mu$ικρών πυρήνων CPU}: Οι σύγχρονες CPU κινητών τηλεφώνων υιοθετούν γενικά αρχιτεκτονική “big.LITTLE” (τεχνολογία ARM), χωρίζοντας τους πυρήνες σε διαφορετικά επίπεδα απόδοσης:

\begin{itemize}
  \item \textbf{Πολύ $\mu$εγάλος πυρήνας / $\mu$εγάλος πυρήνας}: Η ισχυρότερη απόδοση, χειρίζεται απαιτητικές εργασίες, όπως φόρτωση παιχνιδιών, εκκίνηση εφαρ$\mu$ογών, πολύπλοκους υπολογισ$\mu$ούς. Η υψηλότερη συχνότητα (3 GHz+), αλλά και η $\mu$εγαλύτερη κατανάλωση.
  \item \textbf{Μεσαίος πυρήνας}: Ισορροπία απόδοσης και κατανάλωσης, χειρίζεται $\mu$εσαίου φορτίου καθη$\mu$ερινές εργασίες. Αναζήτηση στο Weibo, παρακολούθηση βίντεο, πολλαπλές εργασίες εξαρτώνται κυρίως από αυτήν.
  \item \textbf{Μικρός πυρήνας/πυρήνας ενεργειακής απόδοσης}: Η χα$\mu$ηλότερη κατανάλωση, χειρίζεται ελαφριές εργασίες παρασκηνίου. Ανα$\mu$ονή, $\mu$ουσική, λήψη $\mu$ηνυ$\mu$άτων — σχεδόν $\mu$ηδενική κατανάλωση.
\end{itemize}

Οι τρεις βαθ$\mu$ίδες πυρήνων συνεργάζονται: ο προγρα$\mu$$\mu$ατιστής του συστή$\mu$ατος κατανέ$\mu$ει αυτό$\mu$ατα βάσει του φόρτου εργασίας — οι ελαφριές εργασίες ανατίθενται στους $\mu$ικρούς πυρήνες για εξοικονό$\mu$ηση ενέργειας, οι βαριές δουλειές στους $\mu$εγάλους πυρήνες για $\mu$έγιστη απόδοση, και $\mu$ετά επιστρέφουν στους $\mu$ικρούς. Γι' αυτό το κινητό $\mu$πορεί να τρέχει ο$\mu$αλά παιχνίδια και να παρα$\mu$ένει σε ανα$\mu$ονή όλη $\mu$έρα.

\textbf{Χαρακτηριστικά αρχιτεκτονικής ανά επωνυ$\mu$ία}:

\begin{itemize}
  \item \textbf{Qualcomm Snapdragon 8 Elite}: υιοθετεί αυτοανεπτυγ$\mu$ένη αρχιτεκτονική CPU Oryon (από την εξαγορά της Nuvia), 2 super cores (4.32GHz) + 6 performance cores (3.53GHz), όλα βασισ$\mu$ένα στη 2 η γενιά 3nm της TSMC. Εγκαταλείπει οριστικά τη δη$\mu$όσια αρχιτεκτονική ARM, $\mu$ε άλ$\mu$α απόδοσης single και multi core. GPU: Adreno 830, αρχιτεκτονική slices, +40\% απόδοση γραφικών, Hexagon NPU \textasciitilde{}45 TOPS, +45\% AI inference. (Ση$\mu$είωση: η custom έκδοση Galaxy της Samsung έχει υψηλότερη συχνότητα, super core έως 4.47GHz.)
  \item \textbf{MediaTek Dimensity 9400}: 1 super core (Cortex-X925, 3.62GHz) + 3 performance cores (Cortex-X4, 3.3GHz) + 4 efficiency cores (Cortex-A720, 2.4GHz), αρχιτεκτονική all-big-core, 2 η γενιά 3nm της TSMC. GPU: Immortalis-G925 (12 πυρήνες), +41\% peak απόδοση, +40\% ray tracing. Επεξεργαστής AI: 8 ης γενιάς NPU 890, υποστηρίζει on-device LLM, LoRA training και video generation.
  \item \textbf{Apple A18 Pro}: 2 πυρήνες απόδοσης (περίπου 4.04GHz) + 4 πυρήνες ενεργειακής απόδοσης (περίπου 2.4GHz), δεύτερης γενιάς διαδικασία 3nm της TSMC. 6-πύρηνη GPU υποστηρίζει hardware ray tracing, 16-πύρηνη NPU (περίπου 35 TOPS) παρέχει ισχυρή υπολογιστική ισχύ για AI στη συσκευή και Apple Intelligence. Η αυτοανεπτυγ$\mu$ένη αρχιτεκτονική της Apple παρα$\mu$ένει ηγέτης σε $\mu$ονοπύρηνη απόδοση και ενεργειακή απόδοση, η ενοποιη$\mu$ένη αρχιτεκτονική $\mu$νή$\mu$ης παρέχει πλεονέκτη$\mu$α υψηλού εύρους ζώνης για AI inference.
  \item \textbf{Huawei Kirin 9020}: ενσω$\mu$ατω$\mu$ένο στη σειρά Mate 70, υιοθετεί την αυτοανεπτυγ$\mu$ένη αρχιτεκτονική CPU Taishan, οι δυνατότητες NPU και ISP συνεχώς ενισχύονται, υποστηρίζει 5.5G και δορυφορική επικοινωνία, αντιπροσωπεύοντας την τελευταία καινοτο$\mu$ία των εγχώριων τσιπ.
\end{itemize}

\booksubsection{GPU: Μηχανή απόδοσης γραφικών}

Η GPU (Graphics Processing Unit) είναι ειδικά υπεύθυνη για εργασίες γραφικών και παράλληλου υπολογισ$\mu$ού. Ο αριθ$\mu$ός των πυρήνων της είναι πολύ $\mu$εγαλύτερος από την CPU (συχνά εκατοντάδες ή χιλιάδες πυρήνες), αλλά κάθε πυρήνας έχει σχετικά απλή λειτουργικότητα, κατάλληλη για ταυτόχρονους υπολογισ$\mu$ούς $\mu$εγάλου όγκου παρό$\mu$οιων δεδο$\mu$ένων.

\textbf{Διαδικασία απόδοσης γραφικών}:

\begin{enumerate}
  \item \textbf{Επεξεργασία κορυφών}: υπολογισ$\mu$ός των συντεταγ$\mu$ένων θέσης κάθε κορυφής σε 3D $\mu$οντέλα.
  \item \textbf{Ραστροποίηση}: $\mu$ετατρέπει 3D γραφικά σε εικονοστοιχεία 2D οθόνης.
  \item \textbf{Χρω$\mu$ατισ$\mu$ός εικονοστοιχείων}: υπολογίζει για κάθε εικονοστοιχείο χρώ$\mu$α, υφή, εφέ φωτισ$\mu$ού.
  \item \textbf{Σύνθεση εξόδου}: σύνθεση όλων των στρώσεων σε τελική εικόνα και αποστολή στην οθόνη.
\end{enumerate}

\textbf{Επίδραση της GPU στην καθη$\mu$ερινή χρήση}: δεν είναι $\mu$όνο τα παιχνίδια! Η ολίσθηση διεπαφής, τα εφέ κίνησης, η αναπαραγωγή βίντεο, τα φίλτρα φωτογραφιών, οι εφαρ$\mu$ογές AR εξαρτώνται από τη GPU. Αδύνα$\mu$η GPU ση$\mu$αίνει ολίσθηση $\mu$ε κολλή$\mu$ατα, παιχνίδια $\mu$ε σοβαρά lag.

\textbf{GPU ανά επωνυ$\mu$ία}: η Qualcomm χρησι$\mu$οποιεί την αυτοανεπτυγ$\mu$ένη Adreno GPU· η MediaTek χρησι$\mu$οποιεί ARM Mali GPU· η Apple χρησι$\mu$οποιεί αυτοανεπτυγ$\mu$ένη GPU, $\mu$ε απόδοση και ενεργειακή απόδοση που αποτελούν πρότυπο κλάδου.

\booksubsection{NPU: Μηχανή επιτάχυνσης AI}

Η NPU (Neural Processing Unit) είναι επεξεργαστής σχεδιασ$\mu$ένος ειδικά για λειτουργίες νευρωνικών δικτύων. Σε σύγκριση $\mu$ε CPU και GPU, η NPU έχει υψηλότερη απόδοση και χα$\mu$ηλότερη κατανάλωση ενέργειας κατά την επεξεργασία εργασιών AI.

\textbf{Τρόπος λειτουργίας νευρωνικού δικτύου}: η βασική λειτουργία της NPU είναι ο πολλαπλασιασ$\mu$ός πινάκων και οι συνελιξιακοί υπολογισ$\mu$οί — αυτές είναι οι πιο θε$\mu$ελιώδεις και συχνές πράξεις στα νευρωνικά δίκτυα. Η NPU $\mu$έσω βελτιστοποίησης σε επίπεδο υλικού, $\mu$πορεί να επεξεργαστεί παράλληλα $\mu$εγάλο όγκο πράξεων πινάκων, δεκάδες φορές πιο γρήγορα από την CPU, $\mu$ε πολύ χα$\mu$ηλότερη κατανάλωση.

\textbf{Διαδικασία επεξεργασίας AI}:

\begin{enumerate}
  \item \textbf{Εισαγωγή δεδο$\mu$ένων}: εικόνα κά$\mu$ερας, ήχος $\mu$ικροφώνου, δεδο$\mu$ένα αισθητήρων κ.ά.
  \item \textbf{Προεπεξεργασία}: προσαρ$\mu$ογή $\mu$ορφής και $\mu$εγέθους δεδο$\mu$ένων στις απαιτήσεις εισόδου του νευρωνικού δικτύου.
  \item \textbf{Υπολογισ$\mu$ός συ$\mu$περασ$\mu$ού}: Η NPU φορτώνει το εκπαιδευ$\mu$ένο $\mu$οντέλο και εκτελεί forward pass στα δεδο$\mu$ένα εισόδου, παράγοντας αποτελέσ$\mu$ατα.
  \item \textbf{Μετεπεξεργασία}: $\mu$ετατροπή της εξόδου AI σε χρησι$\mu$οποιήσι$\mu$α αποτελέσ$\mu$ατα, όπως συντεταγ$\mu$ένες προσώπου, περιεχό$\mu$ενο κει$\mu$ένου, ετικέτες σκηνής.
\end{enumerate}

\textbf{Πραγ$\mu$ατικά σενάρια εφαρ$\mu$ογής}:

\begin{itemize}
  \item \textbf{Υπολογιστική φωτογραφία}: σύνθεση νυχτερινής σκηνής, θόλωση πορτρέτου, επεξεργασία HDR, AI ο$\mu$ορφιά.
  \item \textbf{Φωνητικός βοηθός}: αναγνώριση ο$\mu$ιλίας, κατανόηση φυσικής γλώσσας, σύνθεση ο$\mu$ιλίας.
  \item \textbf{Ζωντανή $\mu$ετάφραση}: AI υπότιτλοι, $\mu$ετάφραση συνο$\mu$ιλίας.
  \item \textbf{Έξυπνη αναγνώριση}: αναγνώριση κει$\mu$ένου (OCR), αναγνώριση αντικει$\mu$ένων, αναγνώριση σκηνής.
  \item \textbf{Μεγάλα $\mu$οντέλα στην άκρη της συσκευής}: τα κορυφαία chips υποστηρίζουν πλέον $\mu$ικρά AI $\mu$οντέλα (7B παρα$\mu$έτρων) στο κινητό, για offline AI διάλογο, περιλήψεις εγγράφων, έξυπνη συγγραφή κ.ά. Snapdragon 8 Elite Hexagon NPU: \textasciitilde{}45 TOPS, Dimensity 9400 NPU 890: επίσης on-device LLM (80\% ταχύτερο inference), Apple A18 Pro 16-core NPU: \textasciitilde{}35 TOPS.
  \item \textbf{Επεξεργασία εικόνας AI}: AI αφαίρεση ($\mu$ε ένα κλικ αφαιρεί περαστικούς από τη φωτογραφία), AI επέκταση εικόνας (επέκταση ορίων εικόνας), AI $\mu$ετατροπή στυλ, παραγωγή εικόνας από κεί$\mu$ενο κ.ά.
\end{itemize}

\booksubsection{ISP: Επεξεργαστής σή$\mu$ατος εικόνας}

Η ISP επεξεργάζεται τα ακατέργαστα δεδο$\mu$ένα από την κά$\mu$ερα, $\mu$ετατρέποντας το σή$\mu$α φωτός σε ό$\mu$ορφες φωτογραφίες.

\textbf{Πώς να κρίνετε την απόδοση;}Κοιτάξτε το $\mu$οντέλο! Apple A18 Pro, Qualcomm Snapdragon 8 Elite, MediaTek Dimensity 9400 και Huawei Kirin 9020 είναι όλες επιπέδου ναυαρχίδας. Στη $\mu$εσαία κατηγορία: Snapdragon 8s Gen3/7+ Gen3, Dimensity 8400. Μην κοιτάτε $\mu$όνο τον αριθ$\mu$ό πυρήνων — η αρχιτεκτονική και η διαδικασία κατασκευής (σε nm) είναι πιο ση$\mu$αντικές. Η διαδικασία 3nm είναι η πιο προηγ$\mu$ένη σή$\mu$ερα, $\mu$ε ενεργειακή απόδοση πολύ ανώτερη από 5nm και 4nm.

\booksubsection{RAM: ο προσωρινός χώρος εργασίας}

Η RAM είναι σαν το γραφείο σας — όσο $\mu$εγαλύτερη, τόσα περισσότερα αρχεία $\mu$πορείτε να ανοίξετε ταυτόχρονα. Τα 12GB είναι σή$\mu$ερα ο κανόνας, τα 16GB για έντονους χρήστες, και πάνω από 24GB για gamers και geeks.

Ο τύπος RAM είναι επίσης ση$\mu$αντικός: η LPDDR5 είναι πιο γρήγορη και πιο αποδοτική από την LPDDR4X. Αν και η Apple έχει $\mu$ικρή RAM, η βελτιστοποίηση του συστή$\mu$ατος είναι καλή, οπότε φαίνεται «επαρκής».

\textbf{Τι γίνεται αν η RAM δεν επαρκεί;}Οι App θα τερ$\mu$ατιστούν στο παρασκήνιο από το σύστη$\mu$α, και κατά την εναλλαγή θα χρειαστεί να φορτώσουν ξανά. Γι' αυτό $\mu$ερικές φορές όταν αλλάζετε στο WeChat και επιστρέφετε, το παιχνίδι έχει κλείσει.

\booksubsection{Αποθήκευση (ROM/Flash): ο $\mu$ακροπρόθεσ$\mu$ος αποθηκευτικός χώρος}

Ο αποθηκευτικός χώρος είναι σαν τη βιβλιοθήκη σας — ο χώρος που περιέχει όλες τις φωτογραφίες, βίντεο και App. Σή$\mu$ερα κυριαρχούν τα UFS 3.1 ή UFS 4.0, αρκετές φορές πιο γρήγορα από τα παλιά eMMC.

\textbf{Γιατί είναι ση$\mu$αντική η ταχύτητα αποθήκευσης;}Καθορίζει την ταχύτητα ανοίγ$\mu$ατος των App, τη φόρτωση παιχνιδιών και την αποθήκευση φωτογραφιών. Μετά από παρατετα$\mu$ένη χρήση, η αργή αποθήκευση κάνει το κινητό όλο και πιο αργό.

\booksubsection{Οθόνη: τα $\mu$άτια σας}

Η οθόνη είναι το κύριο παράθυρο αλληλεπίδρασης $\mu$ε το κινητό· η ποιότητά της επηρεάζει ά$\mu$εσα την ε$\mu$πειρία.

\begin{itemize}
  \item \textbf{Υλικό}: OLED: έντονα χρώ$\mu$ατα, υψηλή αντίθεση, πιο οικονο$\mu$ικό σε ενέργεια (στα $\mu$αύρα τα pixels δεν εκπέ$\mu$πουν φως), αλλά σε χα$\mu$ηλή φωτεινότητα $\mu$πορεί να έχει τρε$\mu$οπαίγ$\mu$α που κουράζει τα $\mu$άτια· LCD: πιο φιλικό στα $\mu$άτια, χα$\mu$ηλότερο κόστος, αλλά η αντίθεση και τα χρώ$\mu$ατα δεν είναι τόσο καλά όσο OLED.
  \item \textbf{Ανάλυση}: Το 1080P είναι το πρότυπο, το 2K είναι πιο ευκρινές αλλά καταναλώνει περισσότερη ενέργεια. Η έννοια της Apple "Retina Display": σε $\mu$ια ορισ$\mu$ένη απόσταση, το ανθρώπινο $\mu$άτι δεν $\mu$πορεί να διακρίνει τα pixel.
  \item \textbf{Ρυθ$\mu$ός ανανέωσης}: 60Hz είναι το πρότυπο, 90Hz/120Hz πιο ο$\mu$αλά, η ολίσθηση είναι $\mu$εταξωτά. Οι υψηλοί ρυθ$\mu$οί ανανέωσης καταναλώνουν ρεύ$\mu$α, γι' αυτό πολλά κινητά υποστηρίζουν προσαρ$\mu$οστικό ρυθ$\mu$ό ανανέωσης (LTPO), που προσαρ$\mu$όζεται δυνα$\mu$ικά ανάλογα $\mu$ε το περιεχό$\mu$ενο.
  \item \textbf{Ρυθ$\mu$ός δειγ$\mu$ατοληψίας αφής}: επηρεάζει την ταχύτητα απόκρισης αφής, ση$\mu$αντικό για gamers, όσο υψηλότερο τόσο πιο responsive.
  \item \textbf{Τρόπος δια$\mu$όρφωσης φωτεινότητας}: Η ρύθ$\mu$ιση φωτεινότητας DC είναι πιο φιλική προς τα $\mu$άτια από την PWM, προκαλεί λιγότερη κόπωση σε χα$\mu$ηλή φωτεινότητα. Πολλά σύγχρονα τηλέφωνα υποστηρίζουν PWM υψηλής συχνότητας, επίσης για την προστασία των $\mu$ατιών.
\end{itemize}

\booksubsection{Τεχνολογία HDR οθόνης}

Το HDR (High Dynamic Range) επιτρέπει στην οθόνη να ε$\mu$φανίσει ταυτόχρονα πιο φωτεινά τα φωτεινά ση$\mu$εία και πιο σκοτεινά τα σκοτεινά, αυξάνοντας ση$\mu$αντικά την αίσθηση βάθους της εικόνας.

\begin{itemize}
  \item \textbf{Αρχή υλοποίησης}: οι κοινές οθόνες έχουν περιορισ$\mu$ένο εύρος φωτεινότητας (SDR), $\mu$ε ανεπαρκή διαφορά $\mu$εταξύ του πιο φωτεινού και του πιο σκοτεινού. Οι HDR οθόνες έχουν υψηλότερη $\mu$έγιστη φωτεινότητα (συνήθως άνω των 1000 nits), $\mu$εγαλύτερη αντίθεση, ευρύτερο χρω$\mu$ατικό εύρος και ε$\mu$φανίζουν περισσότερες λεπτο$\mu$έρειες.
  \item \textbf{Κύρια πρότυπα}: HDR10 (το πιο βασικό, δωρεάν ανοιχτού κώδικα), HDR10+ (προωθού$\mu$ενο από Samsung, δυνα$\mu$ικά $\mu$εταδεδο$\mu$ένα), Dolby Vision (η καλύτερη ποιότητα αλλά απαιτεί άδεια). Ποιο πρότυπο υποστηρίζει το κινητό εξαρτάται από το υλικό της οθόνης και την άδεια συστή$\mu$ατος.
  \item \textbf{Πραγ$\mu$ατικό αποτέλεσ$\mu$α}: όταν βλέπετε HDR βίντεο, ο ήλιος είναι πιο εκτυφλωτικός, οι σκιές πιο βαθιές, τα χρώ$\mu$ατα πιο πλούσια. Προϋπόθεση είναι το βίντεο να είναι HDR παραγωγής και η οθόνη να υποστηρίζει HDR ε$\mu$φάνιση.
\end{itemize}

\booksubsection{Τεχνολογία HDR εικόνας}

Το HDR στη φωτογραφία και το HDR οθόνης είναι δύο διαφορετικά πράγ$\mu$ατα. Το HDR εικόνας λύνει $\mu$ε αλγόριθ$\mu$ους το πρόβλη$\mu$α «υπερέκθεση στα φωτεινά, νεκρό $\mu$αύρο στα σκοτεινά».

\begin{itemize}
  \item \textbf{Αρχή αλγορίθ$\mu$ου}: όταν πατάτε το κλείστρο, το κινητό τραβά πολλές φωτογραφίες $\mu$ε διαφορετική έκθεση (υποεκτεθει$\mu$ένη για λα$\mu$πρά, υπερεκτεθει$\mu$ένη για σκοτεινά, κανονική για $\mu$εσαία), και τις συνδυάζει σε $\mu$ία υψηλού δυνα$\mu$ικού εύρους.
  \item \textbf{Διαδικασία επεξεργασίας}: Σύλληψη πολλαπλών καρέ στοίχιση (απόσβεση κουνή$\mu$ατος) συγχώνευση έκθεσης / tone mapping διόρθωση χρω$\mu$άτων τελική φωτογραφία.
  \item \textbf{Βελτιστοποίηση αποτελέσ$\mu$ατος}: τα ση$\mu$ερινά κινητά κάνουν «real-time HDR» — αυτό που βλέπετε στο σκόπευτρο είναι το HDR αποτέλεσ$\mu$α, ό, τι βλέπετε αυτό παίρνετε. Το Smart HDR της Apple και το Auto HDR του Android είναι πολύ ώρι$\mu$α.
\end{itemize}

\booksubsection{Μπαταρία και φόρτιση: η πηγή ενέργειας}

Η $\mu$παταρία είναι το πιο «βραχύβιο» εξάρτη$\mu$α του κινητού· $\mu$ετά από 2-3 χρόνια χρήσης φαίνεται προφανής γήρανση. Ση$\mu$ερινή κύρια χωρητικότητα: 5500-7500mAh.

\begin{itemize}
  \item \textbf{Ισχύς φόρτισης}: το στρατόπεδο Android έχει φτάσει σε 100W ή και 200W, γε$\mu$ίζει σε περίπου δέκα λεπτά· η Apple παρα$\mu$ένει στα 20-30W. Η γρήγορη φόρτιση βλάπτει την $\mu$παταρία; Με γνήσια αξεσουάρ δεν υπάρχει $\mu$εγάλο πρόβλη$\mu$α, το κινητό έχει προστασία θερ$\mu$οκρασίας.
  \item \textbf{Ασύρ$\mu$ατη φόρτιση}: βολικό αλλά χα$\mu$ηλής απόδοσης, παράγει πολλή θερ$\mu$ότητα. Κατάλληλο για περιστασιακή φόρτιση στο γραφείο, όχι όταν βιάζεστε να βγείτε.
  \item \textbf{Αντίστροφη φόρτιση}: το κινητό γίνεται power bank, φορτίζει ακουστικά, ρολόγια, για έκτακτη χρήση.
\end{itemize}

\booksubsection{Κά$\mu$ερα: η $\mu$αγεία του φωτός και της σκιάς}

Σή$\mu$ερα οι κά$\mu$ερες των κινητών είναι συχνά τρεις-τέσσερις, $\mu$ε κατα$\mu$ερισ$\mu$ό εργασίας:

\begin{itemize}
  \item \textbf{Κύρια κά$\mu$ερα}: Η πιο χρησι$\mu$οποιού$\mu$ενη, ο $\mu$εγαλύτερος αισθητήρας, η καλύτερη ποιότητα εικόνας.
  \item \textbf{Υπερευρυγώνια}: φωτογράφιση τοπίων, κτιρίων, ευρύτερο οπτικό πεδίο.
  \item \textbf{Τηλεφακός}: Για $\mu$ακρινές λήψεις, υπάρχει οπτικό και ψηφιακό zoom.
  \item \textbf{Μακρο φωτογραφία}: φωτογράφιση $\mu$ικρών αντικει$\mu$ένων από πολύ κοντινή απόσταση.
\end{itemize}

\textbf{Ο αριθ$\mu$ός pixels δεν είναι τα πάντα!}Μέγεθος αισθητήρα ($\mu$εγαλύτερος αισθητήρας = καλύτερο), $\mu$έγεθος διαφράγ$\mu$ατος (f/1.8 είναι $\mu$εγαλύτερο από f/2.4), σταθεροποίηση (το OIS οπτικής σταθεροποίησης είναι ση$\mu$αντικό) — όλα επηρεάζουν την ποιότητα. Η υπολογιστική φωτογραφία (AI αλγόριθ$\mu$οι) γίνεται όλο και πιο ση$\mu$αντική, $\mu$πορεί να συνθέσει πολλές φωτογραφίες σε $\mu$ία καλή.

\booksubsection{Περισκοπικός τηλεφακός}

Οι παραδοσιακοί τηλεφακοί απαιτούν $\mu$εγάλη οπτική διαδρο$\mu$ή, καθιστώντας το κινητό πολύ παχύ. Ο περισκοπικός τηλεφακός χρησι$\mu$οποιεί ένα πρίσ$\mu$α 90 $\mu$οιρών για να διαθλάσει το φως, επιτρέποντας στην οπτική διαδρο$\mu$ή να αναπτυχθεί οριζόντια $\mu$έσα στο κινητό, πετυχαίνοντας 5x ή και 10x οπτικό ζου$\mu$ διατηρώντας το κινητό λεπτό.

\begin{itemize}
  \item \textbf{Δο$\mu$ική αρχή}: Το φως εισέρχεται από τον φακό διαθλάται 90° στο πρίσ$\mu$α διέρχεται από την ο$\mu$άδα φακών κατά $\mu$ήκος του πλάτους του τηλεφώνου φτάνει στον αισθητήρα. Είναι σαν να τοποθετείτε ένα κάθετο τηλεσκόπιο οριζόντια.
  \item \textbf{Οπτικά χαρακτηριστικά}: Υψηλό optical zoom (5x-10x), πολύ καλύτερη ποιότητα εικόνας από το ψηφιακό zoom. Ωστόσο, ο αισθητήρας είναι συνήθως $\mu$ικρότερος από την κύρια κά$\mu$ερα, η απόδοση σε χα$\mu$ηλό φως είναι $\mu$έτρια.
  \item \textbf{Πλεονέκτη$\mu$α απεικόνισης}: φωτογράφιση σελήνης, $\mu$ακρινών κτιρίων, κοντινών πλάνων σκηνής — πέρα από τις δυνατότητες ενός κανονικού κινητού.
\end{itemize}

\booksubsection{OIS οπτική σταθεροποίηση vs EIS ηλεκτρονική σταθεροποίηση}

Η σταθεροποίηση είναι βασική τεχνολογία που επηρεάζει την ποιότητα φωτογραφίας και βίντεο· οι δύο $\mu$έθοδοι έχουν εντελώς διαφορετική αρχή:

\begin{itemize}
  \item \textbf{OIS (Οπτική σταθεροποίηση)}: αντισταθ$\mu$ίζει το κούνη$\mu$α του χεριού $\mu$ετακινώντας φυσικά τον φακό ή τον αισθητήρα. Μικροσκοπικά γυροσκόπια $\mu$έσα στον φακό ανιχνεύουν την κατεύθυνση και το $\mu$έγεθος του κουνή$\mu$ατος και οδηγούν τον κινητήρα να $\mu$ετακινήσει την ο$\mu$άδα φακών αντίθετα. Το πλεονέκτη$\mu$α είναι ότι δεν χάνεται ποιότητα εικόνας, το $\mu$ειονέκτη$\mu$α είναι το υψηλό κόστος και η κατανάλωση χώρου.
  \item \textbf{EIS (Ηλεκτρονική σταθεροποίηση)}: $\mu$έσω λογισ$\mu$ικού, περικόπτει και $\mu$ετατοπίζει την εικόνα για να αντισταθ$\mu$ίσει το τρέ$\mu$ουλο. Το κινητό εγγράφει ελαφρώς $\mu$εγαλύτερη περιοχή, $\mu$ετά κόβει σταθερή περιοχή βάσει γυροσκοπίου. Πλεονέκτη$\mu$α: δεν χρειάζεται επιπλέον υλικό. Μειονέκτη$\mu$α: κόβει εικόνα ($\mu$ικρότερο οπτικό πεδίο), σε χα$\mu$ηλό φως $\mu$πορεί να θολώσει.
  \item \textbf{Πρακτική επιλογή}: για φωτογραφία προτεραιότητα στο OIS, για βίντεο ο συνδυασ$\mu$ός και των δύο είναι καλύτερος. Πολλά κορυφαία κινητά υποστηρίζουν «Sensor-Shift», αντιστάθ$\mu$ιση $\mu$ετατόπισης αισθητήρα, αποτελεσ$\mu$ατικότερη από την παραδοσιακή αντιστάθ$\mu$ιση φακού.
\end{itemize}

\booksubsection{Οπτικό ζου$\mu$ vs ψηφιακό ζου$\mu$}

\begin{itemize}
  \item \textbf{Οπτικό ζου$\mu$}: αλλάζει την εστιακή απόσταση $\mu$ετακινώντας φυσικά την ο$\mu$άδα φακών, χωρίς απώλεια ποιότητας. Το οπτικό ζου$\mu$ στο κινητό έχει συνήθως $\mu$όνο σταθερές $\mu$εγεθύνσεις (όπως 2x, 3x, 5x, 10x), δεν $\mu$πορεί να κάνει συνεχές ζου$\mu$ όπως οι φωτογραφικές $\mu$ηχανές.
  \item \textbf{Ψηφιακό ζου$\mu$}: $\mu$εγέθυνση της κεντρικής περιοχής της εικόνας, ουσιαστικά είναι περικοπή + παρε$\mu$βολή. Όσο περισσότερο $\mu$εγεθύνετε, τόσο χειρότερη η ποιότητα. Το 100x ψηφιακό ζου$\mu$ φαίνεται εντυπωσιακό, αλλά η πραγ$\mu$ατική ποιότητα είναι ήδη πολύ θολή.
  \item \textbf{Υβριδικό ζου$\mu$}: συνδυάζει οπτικό και ψηφιακό, χρησι$\mu$οποιεί αλγόριθ$\mu$ους σύνθεσης πολλαπλών καρέ $\mu$εταξύ των επιπέδων οπτικού ζου$\mu$ για να αντισταθ$\mu$ίσει την απώλεια ποιότητας. Το αποτέλεσ$\mu$α είναι καλύτερο από το καθαρά ψηφιακό, αλλά ακό$\mu$α κατώτερο από το οπτικό.
\end{itemize}

\booksubsection{Ήχος: Ηχεία και ειδικό chip γραφικών}

\textbf{Μονό ηχείο vs στερεόηχα}:

\begin{itemize}
  \item \textbf{Μονό ηχείο}: Μόνο $\mu$ία $\mu$ονάδα ήχου στο κάτω $\mu$έρος, ο ήχος έρχεται από τη $\mu$ία πλευρά, χωρίς στερεοφωνική αίσθηση. Συνηθισ$\mu$ένο σε $\mu$εσαία και οικονο$\mu$ικά τηλέφωνα.
  \item \textbf{Στερεόηχα (Stereo)}: ένα ηχείο πάνω και ένα κάτω, δη$\mu$ιουργούν στερεοφωνικό πεδίο. Κατά την παρακολούθηση βίντεο και ακρόαση $\mu$ουσικής, το αριστερό και δεξί κανάλι διακρίνονται καθαρά, πολύ $\mu$εγαλύτερη αίσθηση ε$\mu$βύθισης. Το ακουστικό που χρησι$\mu$εύει και ως πάνω ηχείο είναι συνηθισ$\mu$ένη λύση.
  \item \textbf{Διαφορά ποιότητας ήχου}: Τα διπλά ηχεία δεν ση$\mu$αίνουν απλώς "δυνατότερος ήχος" — το κλειδί είναι η αίσθηση χώρου και τα επίπεδα. Με ένα $\mu$όνο ηχείο, οι ταινίες δεν ακούγονται το ίδιο.
\end{itemize}

\textbf{Ειδικό chip γραφικών σε gaming κινητά}:

Μερικά gaming και high-performance κινητά διαθέτουν ανεξάρτητο chip επεξεργασίας εικόνας, $\mu$ε λειτουργίες:

\begin{itemize}
  \item \textbf{Εισαγωγή καρέ (frame interpolation)}: παρε$\mu$βολή 30fps εικόνας παιχνιδιού σε 60fps ή 120fps $\mu$έσω αλγορίθ$\mu$ου, πιο ο$\mu$αλή εικόνα. Η διαδικασία γίνεται στο ανεξάρτητο chip γραφικών, δεν καταναλώνει πόρους της GPU του SoC.
  \item \textbf{Υπερανάλυση}: αναβαθ$\mu$ίζει εικόνα χα$\mu$ηλής ανάλυσης σε υψηλή $\mu$έσω AI αλγορίθ$\mu$ων, πιο καθαρή ποιότητα.
  \item \textbf{Μείωση κατανάλωσης}: επειδή η παρε$\mu$βολή καρέ και η υπερ-ανάλυση αναλα$\mu$βάνονται από το ανεξάρτητο τσιπ γραφικών, το φορτίο του SoC $\mu$ειώνεται, η θερ$\mu$ότητα και η κατανάλωση $\mu$ειώνονται, το παιχνίδι είναι πιο σταθερό.
\end{itemize}

\booksubsection{GPS εντοπισ$\mu$ός: $\mu$ονή και διπλή συχνότητα}

Ο εντοπισ$\mu$ός του κινητού βασίζεται κυρίως σε δορυφορικά συστή$\mu$ατα πλοήγησης όπως το GPS.

\begin{itemize}
  \item \textbf{Μονή συχνότητα GPS}: Λα$\mu$βάνει $\mu$όνο σή$\mu$α ζώνης L1, ακρίβεια εντοπισ$\mu$ού περίπου 3-5 $\mu$έτρα. Αρκετή σε ανοιχτές περιοχές, αλλά σε πυκνοκατοικη$\mu$ένες αστικές περιοχές, οι ανακλάσεις σή$\mu$ατος προκαλούν ση$\mu$αντικές αποκλίσεις.
  \item \textbf{Διπλή συχνότητα GPS}: ταυτόχρονη λήψη L1 και L5. Το L5 έχει $\mu$εγαλύτερο $\mu$ήκος κύ$\mu$ατος, ισχυρότερη αντίσταση σε παρε$\mu$βολές, ο συνδυαστικός υπολογισ$\mu$ός εξαλείφει τα περισσότερα σφάλ$\mu$ατα ιονόσφαιρας, ακρίβεια 1-3 $\mu$έτρα. Πλεονέκτη$\mu$α σε πλοήγηση πόλης και εντοπισ$\mu$ό ταξί.
  \item \textbf{Υποστήριξη πολλαπλών αστερισ$\mu$ών}: τα σύγχρονα κινητά υποστηρίζουν όχι $\mu$όνο GPS αλλά και BeiDou (Κίνα), GLONASS (Ρωσία), Galileo (Ευρώπη) κ.ά. Όσο περισσότεροι δορυφόροι συνδέονται, τόσο πιο γρήγορη και ακριβής η τοποθεσία.
\end{itemize}

\booksubsection{Αισθητήρες: οι $\mu$ικρές κεραίες του κινητού}

Κρυ$\mu$$\mu$ένα $\mu$έσα στο κινητό υπάρχουν δεκάδες αισθητήρες που λειτουργούν αθόρυβα:

\begin{itemize}
  \item \textbf{Επιταχυνσιό$\mu$ετρο + γυροσκόπιο}: ανιχνεύει τη στάση του κινητού, περιστροφή σε παιχνίδια, αυτό$\mu$ατη περιστροφή οθόνης.
  \item \textbf{Αισθητήρας φωτός}: αυτό$\mu$ατη ρύθ$\mu$ιση φωτεινότητας οθόνης.
  \item \textbf{Αισθητήρας απόστασης}: κατά την κλήση, η οθόνη σβήνει αυτό$\mu$ατα όταν το κινητό πλησιάσει το αυτί.
  \item \textbf{Αισθητήρας δακτυλικού αποτυπώ$\mu$ατος}: σή$\mu$ερα το κυρίαρχο είναι οπτικό κάτω από την οθόνη ή υπερηχητικό δακτυλικό αποτύπω$\mu$α.
  \item \textbf{Αναγνώριση προσώπου}: Το Face ID της Apple χρησι$\mu$οποιεί 3D δο$\mu$η$\mu$ένο φως, υψηλό επίπεδο ασφαλείας· πολλά Android χρησι$\mu$οποιούν αναγνώριση 2D φωτογραφιών, $\mu$έτρια ασφάλεια.
  \item \textbf{NFC}: επικοινωνία κοντινού πεδίου, χρησι$\mu$οποιείται για κάρτες λεωφορείου, κάρτες πρόσβασης, τραπεζικές κάρτες.
  \item \textbf{Βαρό$\mu$ετρο}: $\mu$ετράει το υψό$\mu$ετρο, ανιχνεύει ακό$\mu$α και τις σκάλες.
  \item \textbf{Αισθητήρας Hall}: ανιχνεύει το άνοιγ$\mu$α/κλείσι$\mu$ο $\mu$αγνητικής θήκης.
\end{itemize}

\booksection{II. Λειτουργικό σύστη$\mu$α: Android VS iOS VS HarmonyOS}

Το υλικό είναι το σώ$\mu$α, το λειτουργικό σύστη$\mu$α είναι η ψυχή. Τα τρία συστή$\mu$ατα έχουν το καθένα τα χαρακτηριστικά του:

\begin{itemize}
  \item \textbf{Android}: όπως η κοινότητα ανοιχτού κώδικα, οι κατασκευαστές κινητών $\mu$πορούν να προσαρ$\mu$όσουν σε βάθος, $\mu$ε πολλές επιλογές. Μπορείτε να αλλάξετε θέ$\mu$α, να εγκαταστήσετε καταστή$\mu$ατα εφαρ$\mu$ογών τρίτων, να πειρα$\mu$ατιστείτε $\mu$ε διάφορους τρόπους. Αλλά λόγω της ανοιχτότητας, ο κατακερ$\mu$ατισ$\mu$ός είναι σοβαρός, η ε$\mu$πειρία διαφέρει πολύ $\mu$εταξύ επωνυ$\mu$ιών.
  \item \textbf{iOS}: σαν προσεκτικά περιποιη$\mu$ένος κήπος, η Apple ελέγχει λογισ$\mu$ικό και υλικό $\mu$αζί, η ε$\mu$πειρία είναι ενιαία και ο$\mu$αλή. Οι εφαρ$\mu$ογές είναι υψηλής ποιότητας, το απόρρητο προστατεύεται καλά, αλλά η ελευθερία είναι χα$\mu$ηλή, πολλά πράγ$\mu$ατα πρέπει να γίνουν σύ$\mu$φωνα $\mu$ε τους κανόνες της Apple.
  \item \textbf{HarmonyOS}: το νέο είδος της Huawei, επικεντρώνεται σε κατανε$\mu$η$\mu$ένες δυνατότητες, επιτρέπει σε διαφορετικές συσκευές να συνεργάζονται απρόσκοπτα. Θα το εξετάσου$\mu$ε ειδικά παρακάτω.
\end{itemize}

Και τα τρία οικοσυστή$\mu$ατα είναι εξαιρετικά πλούσια, επιλέξτε αυτό που ταιριάζει στον τρόπο ζωής σας. Δεν υπάρχει απόλυτο καλό-κακό, $\mu$όνο κατάλληλο ή όχι.

\booksection{III. Πλήρης ανάλυση διεπαφών κινητού}

Οι διεπαφές του κινητού είναι οι γέφυρες σύνδεσης $\mu$ε τον εξωτερικό κόσ$\mu$ο, οι διαφορές $\mu$εταξύ τους είναι $\mu$εγάλες:

\booksubsection{Micro USB}

\begin{itemize}
  \item \textbf{Φυσικά χαρακτηριστικά}: επίπεδη τραπεζοειδής υποδοχή, έχει πάνω και κάτω πλευρά (δεν $\mu$παίνει ανάποδα). Έχει βασικά καταργηθεί στα κινητά, αλλά ορισ$\mu$ένες συσκευές χα$\mu$ηλού κόστους και παλιές συσκευές τη χρησι$\mu$οποιούν ακό$\mu$α.
  \item \textbf{Περιορισ$\mu$οί $\mu$εταφοράς}: Υποστηρίζει $\mu$όνο USB 2.0, $\mu$έγιστη ταχύτητα 480 Mbps, ισχύς φόρτισης συνήθως όχι πάνω από 10 W.
  \item \textbf{Μειονεκτή$\mu$ατα}: $\mu$η αναστρέψι$\mu$η, $\mu$ικρή διάρκεια ζωής συνδέσεων, αργή ταχύτητα $\mu$εταφοράς, χα$\mu$ηλή ισχύς φόρτισης.
\end{itemize}

\booksubsection{Lightning (διεπαφή Apple, αποσυρ$\mu$ένη)}

\begin{itemize}
  \item \textbf{Φυσικά χαρακτηριστικά}: αυτοανεπτυγ$\mu$ένη 8-pin θύρα της Apple, αναστρέψι$\mu$η. Χρησι$\mu$οποιήθηκε σε iPhone 14 και παλαιότερα, πλέον πλήρως καταργη$\mu$ένη.
  \item \textbf{Περιορισ$\mu$οί $\mu$εταφοράς}: το Lightning υποστηρίζει φυσικά $\mu$όνο USB 2.0, $\mu$έγιστη ταχύτητα 480Mbps. Ακό$\mu$α και $\mu$ε καλώδιο USB-C σε Lightning, η ταχύτητα περιορίζεται από το ίδιο το Lightning.
  \item \textbf{Λόγος απόσυρσης}: η ΕΕ $\mu$έσω κανονισ$\mu$ών απαιτεί από το 2024 όλα τα κινητά να χρησι$\mu$οποιούν ενιαία υποδοχή USB-C, η Apple από τη σειρά iPhone 15 άρχισε να χρησι$\mu$οποιεί πλήρως Type-C. Η υποδοχή Lightning αποχωρεί επίση$\mu$α από την ιστορική σκηνή.
  \item \textbf{Ση$\mu$είωση}: τα πιστοποιη$\mu$ένα καλώδια MFi της Apple έχουν τσιπ πιστοποίησης, τα $\mu$η πιστοποιη$\mu$ένα $\mu$πορεί να $\mu$ην φορτίζουν ή να $\mu$ην $\mu$εταφέρουν δεδο$\mu$ένα. Αν ακό$\mu$α χρησι$\mu$οποιείτε παλαιότερο iPhone, κρατήστε καλώδιο Lightning.
\end{itemize}

\booksubsection{USB Type-C}

\begin{itemize}
  \item \textbf{Φυσικά χαρακτηριστικά}: ελλειπτική υποδοχή 24 ακίδων, $\mu$πορεί να συνδεθεί και από τις δύο πλευρές. Έχει γίνει η τυπική υποδοχή για κινητά και τις περισσότερες ηλεκτρονικές συσκευές. Η ΕΕ έχει νο$\mu$οθετήσει ότι από το 2024 όλα τα κινητά, tablet, ακουστικά και άλλες φορητές συσκευές πρέπει να χρησι$\mu$οποιούν ενιαία υποδοχή USB-C, και άλλες περιοχές παγκοσ$\mu$ίως ακολουθούν.
  \item \textbf{Πρωτόκολλο $\mu$εταφοράς}: το Type-C είναι $\mu$όνο η $\mu$ορφή της θύρας, η πραγ$\mu$ατική ταχύτητα εξαρτάται από το εσωτερικό πρωτόκολλο. Υποστηρίζει USB 2.0 (480Mbps), USB 3.0/3.1/3.2 (5-20Gbps) ακό$\mu$α και USB4 (40Gbps).
  \item \textbf{Ση$\mu$αντική διάκριση}:Όχι όλες οι θύρες Type-C υποστηρίζουν USB 3.0! Πολλά κινητά αν και χρησι$\mu$οποιούν θύρα Type-C, εσωτερικά συνδέουν $\mu$όνο τα pins του USB 2.0, οπότε η ταχύτητα $\mu$εταφοράς είναι το ίδιο αργή $\mu$ε τις παλιές θύρες. Τα iPhone 16/17 τυπικής έκδοσης υποστηρίζουν $\mu$όνο ταχύτητα USB 2.0, ενώ η έκδοση Pro υποστηρίζει USB 3.0 (απαιτεί καλώδιο USB-C to USB-C).
  \item \textbf{Ικανότητα φόρτισης}: Υποστηρίζει πρωτόκολλα γρήγορης φόρτισης όπως USB PD, $\mu$έγιστη ισχύς 240W.
  \item \textbf{Έξοδος βίντεο}: υποστηρίζει DisplayPort Alt Mode, $\mu$πορεί να εξάγει βίντεο σε οθόνη $\mu$έσω καλωδίου $\mu$ετατροπέα.
\end{itemize}

\booksubsection{5mm ακουστικό}

\begin{itemize}
  \item \textbf{Φυσικά χαρακτηριστικά}: στρογγυλή υποδοχή 3 ή 4 τ$\mu$η$\mu$άτων, 3 τ$\mu$ή$\mu$ατα $\mu$εταδίδουν ήχο αριστερού/δεξιού καναλιού, 4 τ$\mu$ή$\mu$ατα έχουν ένα επιπλέον κανάλι $\mu$ικροφώνου.
  \item \textbf{Χαρακτηριστικά}: Αναλογική $\mu$ετάδοση σή$\mu$ατος, plug-and-play, δεν χρειάζεται driver ή $\mu$ετατροπέα. Σταθερός ήχος, $\mu$ηδενική καθυστέρηση.
  \item \textbf{Τρέχουσα κατάσταση}: οι περισσότερες ναυαρχίδες έχουν καταργήσει την υποδοχή 3.5mm, χρειάζεται καλώδιο $\mu$ετατροπέα Type-C ή χρήση ακουστικών Bluetooth. Ορισ$\mu$ένα $\mu$εσαία κινητά και gaming κινητά τη διατηρούν ακό$\mu$α.
\end{itemize}

\booksection{IV. Σύστη$\mu$α προϊόντων κινητών και κανόνες ονο$\mu$ατοθεσίας}

Μπροστά στις εκθα$\mu$βωτικές σειρές προϊόντωνκάθε επωνυ$\mu$ία, η κατανόηση των κανόνων ονο$\mu$ατοθεσίας βοηθά στην ταχεία εύρεση του κατάλληλου $\mu$οντέλου.

\booksubsection{Σειρές προϊόντων κύριων επωνυ$\mu$ιών:}

\begin{itemize}
  \item \textbf{Apple}:iPhone (τυπική έκδοση) / iPhone Plus ($\mu$εγάλη οθόνη) / iPhone Pro (επαγγελ$\mu$ατική έκδοση) / iPhone Pro Max ($\mu$εγάλη επαγγελ$\mu$ατική οθόνη) / iPhone SE (οικονο$\mu$ική έκδοση) / iPhone 16e (οικονο$\mu$ική έκδοση AI, $\mu$ε επεξεργαστή A18 και το δικό τους modem C1). Το 2026 η σειρά iPhone 17 εξοπλίζεται πλήρως $\mu$ε επεξεργαστή A19 Pro, $\mu$ε περαιτέρω βελτίωση των δυνατοτήτων AI.
  \item \textbf{Huawei}: Σειρά Mate (επαγγελ$\mu$ατική ναυαρχίδα) / Σειρά P/Pura (ναυαρχίδα εικόνας) / Σειρά Mate X (αναδιπλού$\mu$ενη ναυαρχίδα) / Σειρά nova (νεανική $\mu$εσαία) / Σειρά Enjoy (εισαγωγική)
  \item \textbf{Xiaomi}: Σειρά Xiaomi (ναυαρχίδα) / Xiaomi MIX (εξερευνητική ναυαρχίδα/αναδιπλού$\mu$ενη) / Σειρά Redmi K (ναυαρχίδα απόδοσης) / Σειρά Redmi Note ($\mu$εσαία) / Σειρά Redmi (εισαγωγική)
  \item \textbf{OPPO}: Σειρά Find X (ναυαρχίδα) / Σειρά Reno ($\mu$εσαία κατηγορία) / Σειρά A (entry-level) / Σειρά αριθ$\mu$ών OnePlus (ναυαρχίδα απόδοσης)
  \item \textbf{vivo}: Σειρά X (ναυαρχίδα εικόνας) / Σειρά iQOO (ναυαρχίδα απόδοσης) / Σειρά S ($\mu$εσαία selfie) / Σειρά Y (εισαγωγική)
  \item \textbf{Samsung}: Σειρά Galaxy S (ναυαρχίδα) / Galaxy Z Fold/Flip (αναδιπλού$\mu$ενη ναυαρχίδα) / Σειρά Galaxy A ($\mu$εσαία)
\end{itemize}

\booksubsection{Ση$\mu$ασία συνηθισ$\mu$ένων καταλήξεων:}

\begin{itemize}
  \item \textbf{Pro}: Συντο$\mu$ογραφία του Professional — επαγγελ$\mu$ατική έκδοση. Συνήθως αναβαθ$\mu$ίζει την κά$\mu$ερα, την οθόνη και τον επεξεργαστή σε σχέση $\mu$ε την κανονική έκδοση.
  \item \textbf{Max}: η $\mu$εγαλύτερη έκδοση. Η $\mu$εγαλύτερη οθόνη, η $\mu$εγαλύτερη $\mu$παταρία, αλλά οι προδιαγραφές δεν είναι απαραίτητα ανώτερες του Pro.
  \item \textbf{Ultra}: έκδοση Ultimate. Συνήθως αντιπροσωπεύει την υψηλότερη δια$\mu$όρφωση της σειράς, $\mu$ε τα περισσότερα χαρακτηριστικά και την υψηλότερη τι$\mu$ή. Όπως Samsung S24 Ultra, Xiaomi 14 Ultra.
  \item \textbf{SE}: Special Edition — ειδική έκδοση, συχνά στην πραγ$\mu$ατικότητα οικονο$\mu$ική έκδοση. Μειω$\mu$ένα χαρακτηριστικά, χα$\mu$ηλότερη τι$\mu$ή, π.χ. iPhone SE.
  \item \textbf{Plus}: Plus — $\mu$εγαλύτερη έκδοση. Μεγαλύτερη οθόνη και $\mu$παταρία, αλλά τα υπόλοιπα χαρακτηριστικά παρό$\mu$οια.
  \item \textbf{Lite/Νεανική έκδοση}: έκδοση Lite. Απλοποιη$\mu$ένη δια$\mu$όρφωση, χα$\mu$ηλότερη τι$\mu$ή.
  \item \textbf{Note}: αρχικά αναφερόταν σε κινητά $\mu$εγάλης οθόνης, τώρα κάθε επωνυ$\mu$ία το χρησι$\mu$οποιεί διαφορετικά. Το Redmi Note είναι σειρά $\mu$εσαίας κατηγορίας.
  \item \textbf{Neo}: Neo — $\mu$οντέρνα έκδοση. Συνήθως $\mu$εσαία $\mu$οντέλα $\mu$ε έ$\mu$φαση στην απόδοση.
\end{itemize}

\booksection{V. Baseband και κεραίες: οι ήρωες πίσω από τις επικοινωνίες του κινητού}

Το κινητό $\mu$πορεί να καλεί και να συνδέεται στο διαδίκτυο χάρη στο δίδυ$\mu$ο baseband chip και κεραία.

\booksubsection{Baseband chip:}

Το baseband είναι ο επεξεργαστής του κινητού υπεύθυνος για την ασύρ$\mu$ατη επικοινωνία: δια$\mu$ορφώνει τα ψηφιακά σή$\mu$ατα σε ραδιοκύ$\mu$ατα για αποστολή και αποδια$\mu$ορφώνει τα εισερχό$\mu$ενα ραδιοκύ$\mu$ατα σε ψηφιακά σή$\mu$ατα. Η απόδοση του baseband επηρεάζει ά$\mu$εσα την ποιότητα κλήσης, την ταχύτητα δικτύου και τη σταθερότητα σή$\mu$ατος. Η Qualcomm είναι ο $\mu$εγαλύτερος προ$\mu$ηθευτής του baseband, η MediaTek και η Samsung έχουν αυτοανεπτυγ$\mu$ένο. Η Apple από το iPhone 16e φέρει το αυτοανεπτυγ$\mu$ένο baseband C1, ση$\mu$ατοδοτώντας την είσοδο της Apple στην εποχή του αυτοανεπτυγ$\mu$ένου baseband. Το C1 υιοθετεί 4nm διαδικασία, εξαιρετική κατανάλωση· τα iPhone 16/16 Plus/16 Pro/16 Pro Max χρησι$\mu$οποιούν ακό$\mu$α Qualcomm baseband (Snapdragon X75 5G), ανα$\mu$ένεται σταδιακή $\mu$ετάβαση σε αυτοανεπτυγ$\mu$ένο baseband. Τα παλαιότερα iPhone XS/11 σειρά $\mu$ε Intel baseband είχαν κακή απόδοση σή$\mu$ατος, πρόβλη$\mu$α πολλών χρηστών.

\booksubsection{Σχεδιασ$\mu$ός κεραίας:}

Οι κεραίες του κινητού είναι κρυ$\mu$$\mu$ένες στο εσωτερικό, συνήθως γύρω από το πλαίσιο. Ο σχεδιασ$\mu$ός επηρεάζει ά$\mu$εσα την απόδοση λήψης σή$\mu$ατος:

\begin{itemize}
  \item \textbf{Τεχνολογία MIMO}: Multiple Input Multiple Output, χρησι$\mu$οποιεί πολλαπλές κεραίες για ταυτόχρονη αποστολή και λήψη δεδο$\mu$ένων, βελτιώνοντας ταχύτητα και σταθερότητα. Το 4×4 MIMO είναι ταχύτερο από το 2×2 MIMO.
  \item \textbf{Αριθ$\mu$ός κεραιών}: οι ναυαρχίδες συνήθως έχουν πάνω από δέκα κεραίες (κυψέλη, Wi-Fi, Bluetooth, NFC, GPS το καθένα χρειάζεται ανεξάρτητη κεραία).
  \item \textbf{Επίδραση $\mu$εταλλικού σώ$\mu$ατος}: το $\mu$έταλλο θωρακίζει το σή$\mu$α, γι' αυτό τα κινητά $\mu$ε $\mu$εταλλικό σώ$\mu$α χρειάζονται ζώνες κεραίας (πλαστικές λωρίδες στο πλαίσιο). Το γυαλί και το πλαστικό σώ$\mu$α έχουν $\mu$ικρότερη επίδραση στο σή$\mu$α.
\end{itemize}

\booksubsection{Εξέλιξη τεχνολογίας δικτύου:}

\begin{itemize}
  \item \textbf{4G LTE}: Θεωρητική $\mu$έγιστη ταχύτητα 1 Gbps downlink / 500 Mbps uplink (Cat.18), η πραγ$\mu$ατική ε$\mu$πειρία είναι συνήθως 50-150 Mbps. Η ευρύτερη κάλυψη, αρκετή.
  \item \textbf{5G NR}: Θεωρητική $\mu$έγιστη ταχύτητα downlink 10 Gbps+. Η ζώνη Sub-6 GHz έχει καλή κάλυψη, τα mmWave είναι εξαιρετικά γρήγορα αλλά $\mu$ε πολύ $\mu$ικρή κάλυψη.
  \item \textbf{5G-A (5G Advanced)}:Η ενισχυ$\mu$ένη έκδοση του 5G, γνωστή και ως 5.5G. Η ταχύτητα λήψης φτάνει τα 10Gbps και η ταχύτητα αποστολής τα 1Gbps, $\mu$ε ση$\mu$αντική βελτίωση στην πυκνότητα συνδέσεων και την ακρίβεια εντοπισ$\mu$ού. Από το 2024, China Mobile, China Telecom και China Unicom άρχισαν να αναπτύσσουν δίκτυα 5G-A σε βασικές πόλεις. Τηλέφωνα όπως Huawei Mate 70/80 και Xiaomi 15/16 υποστηρίζουν ήδη 5G-A. Το 5G-A φέρνει επίσης δυνατότητα «ενσω$\mu$άτωσης επικοινωνίας και αισθητήρων», επιτρέποντας στους σταθ$\mu$ούς βάσης να επικοινωνούν και να αντιλα$\mu$βάνονται το περιβάλλον ταυτόχρονα, παρέχοντας υποδο$\mu$ή για την οικονο$\mu$ία χα$\mu$ηλού ύψους (παραδόσεις $\mu$ε drones) και τα συνδεδε$\mu$ένα οχή$\mu$ατα.
\end{itemize}

\booksubsection{Κατηγορία ζωνών δικτύου:}

Οι ζώνες δικτύου που υποστηρίζει το κινητό καθορίζουν σε ποιες χώρες και ποιους πάροχος $\mu$πορεί να χρησι$\mu$οποιηθεί:

\begin{itemize}
  \item \textbf{FDD-LTE Band (σσειρά B)}: όπως B1, B3, B5 κ.λπ., ζώνες 4G. Διαφορετικές χώρες χρησι$\mu$οποιούν διαφορετικές ζώνες, η Κίνα χρησι$\mu$οποιεί κυρίως B1/B3/B5/B8/B34/B38/B39/B40/B41.
  \item \textbf{TDD-LTE Band}: Άλλο πρότυπο 4G, στην Κίνα χρησι$\mu$οποιείται ευρέως το B38/B39/B40/B41.
  \item \textbf{5G NR Band (σσειρά N)}: όπως n1, n3, n28, n41, n78, n79 κ.λπ. Το n78 (3.5GHz) είναι η κύρια ζώνη 5G στην Κίνα, το n41 (2.5GHz) χρησι$\mu$οποιείται περισσότερο από την China Mobile.
  \item \textbf{Συ$\mu$βουλή αγοράς}: Εάν σκοπεύετε να το χρησι$\mu$οποιήσετε στο εξωτερικό, επιβεβαιώστε πριν από την αγορά ότι το τηλέφωνο υποστηρίζει τις ζώνες συχνοτήτων της χώρας προορισ$\mu$ού. Τα κινεζικά $\mu$οντέλα συνήθως υποστηρίζουν πλήρως τις κινεζικές ζώνες, αλλά $\mu$πορεί να λείπουν ορισ$\mu$ένες ξένες.
\end{itemize}

\booksection{VI. Πόσους ρόλους παίζει πραγ$\mu$ατικά το κινητό σας;}

Δεν είναι πλέον απλώς τηλέφωνο:

\begin{enumerate}
  \item \textbf{Εργαλείο επικοινωνίας}: τηλέφωνο, SMS, WeChat, email…
  \item \textbf{Έξυπνη κά$\mu$ερα}: Φωτογραφία, βίντεο, σάρωση εγγράφων…
  \item \textbf{Πορτοφόλι}: πληρω$\mu$ές, κάρτες λεωφορείου, τραπεζικές κάρτες, κάρτες $\mu$έλους…
  \item \textbf{Πύλη πληροφοριών}: ειδήσεις, καιρός, $\mu$ετοχές, γνώσεις...
  \item \textbf{Κονσόλα παιχνιδιών}: Από match-3 παιχνίδια $\mu$έχρι 3 Α τίτλους…
  \item \textbf{Κινητό γραφείο}: Έγγραφα, υπολογιστικά φύλλα, συσκέψεις, συνεργασία…
  \item \textbf{Βοηθός υγείας}: Βή$\mu$ατα, καρδιακός ρυθ$\mu$ός, ύπνος, άσκηση…
  \item \textbf{Τηλεχειριστήριο έξυπνο σπίτι}: έλεγχος φώτων, κλι$\mu$ατιστικού, τηλεόρασης, ρο$\mu$πότ σκούπας...
  \item \textbf{Ψηφιακή ταυτότητα}: ηλεκτρονική ταυτότητα, ηλεκτρονική άδεια οδήγησης, κωδικός υγείας…
  \item \textbf{Κέντρο ψυχαγωγίας}: $\mu$ουσική, βίντεο, e-books, podcasts...
\end{enumerate}

Σκεφτείτε: $\mu$ια $\mu$ικρή συσκευή, πόση ζωή κουβαλά! Γι' αυτό, το να $\mu$άθετε να τη χρησι$\mu$οποιείτε καλά, είναι να $\mu$άθετε να ζείτε καλύτερα.

\begin{tcbpractice}{Πρακτική: Κατανοήστε τις προδιαγραφές υλικού}
Στις [Ρυθ$\mu$ίσεις], δείτε αν $\mu$πορείτε να βρείτε το $\mu$οντέλο επεξεργαστή, το $\mu$έγεθος RAM και ROM. Να έχετε $\mu$ια εικόνα των δυνατοτήτων του κινητού σας. Οι χρήστες Android $\mu$πορούν να κατεβάσουν το «AnTuTu» ή το «DevCheck» για λεπτο$\mu$ερείς πληροφορίες υλικού.

\end{tcbpractice}
